\sectionkicker{Source-backed technical notes}
\subsection*{規模が設計原理になる――Scaling LawsからGPT-3へ: Technical Notes}
本欄は記事本文で圧縮した一次資料上の情報を、比較・再検証しやすい形で整理したものである。時系列、確認済みの主張、留意点、一次資料URLを掲載する。
\smallskip
\noindent{\small\textit{Technical Notesでは、一次資料から確認した公開・提供・研究上の事実を記載する。数値・能力評価や提供元・著者による比較表現は、独立再現済みの結果ではなく帰属claimとして扱う。}}\par
\medskip
\subsection*{Theme at a glance}
\addcontentsline{toc}{subsection}{Theme at a glance}
\begin{center}
\footnotesize
\begin{tabularx}{\linewidth}{@{}p{0.20\linewidth}p{0.10\linewidth}p{0.14\linewidth}X@{}}
\toprule
資料 & 位置づけ & 種別 & 時系列 \\
\midrule
Scaling Laws for Neural Language Models & 主要資料 & 論文 & — \\
GPT-3 & 主要資料 & モデル & — \\
\bottomrule
\end{tabularx}
\end{center}
\normalsize

\begin{technicalnote}{Scaling Laws for Neural Language Models}{主要資料}
\begingroup
% reader-facing Technical Notes break policy
\widowpenalty=10000
\clubpenalty=10000
\displaywidowpenalty=10000
\begin{tabularx}{\linewidth}{@{}>{\bfseries}p{0.22\linewidth}X@{}}
組織 & OpenAI / authors \\
種別 & 論文 \\
時系列 & — \\
\end{tabularx}
\smallskip
{\bfseries 一次資料から整理したtechnical points}
\begin{itemize}[leftmargin=1.5em,itemsep=0.35em]
\item \textbf{一次情報で確認できる事実}: OpenAI / authorsの選定済み一次資料では、Scaling Laws論文はcross-entropy lossがmodel size・dataset size・training computeに対してpower-lawで変化し、7桁超の範囲でその傾向を観測したと報告する。 固定compute budgetの配分則を導き、著者らは大規模modelほどsample-efficientで、compute-efficient trainingでは非常に大きなmodelを比較的少ないdataで学習し収束前に停止する構成になると述べる。 network widthやdepthなど個別architecture詳細の効果が広い範囲で小さいという記述は、この論文の実験条件における著者観測として扱う。
\end{itemize}
\begin{minipage}{\linewidth}
% reader-facing Technical Notes generic boundary/source tail group
\begin{samepage}
{\bfseries 一次資料}
\begingroup\sloppy
\Urlmuskip=0mu plus 2mu\relax
\begin{itemize}[leftmargin=1.5em,itemsep=0.25em]
\item {\scriptsize\url{http://arxiv.org/abs/2001.08361v1}}
\end{itemize}
\endgroup
\end{samepage}
\end{minipage}
\endgroup
\end{technicalnote}
\medskip

\begin{technicalnote}{GPT-3}{主要資料}
\begingroup
% reader-facing Technical Notes break policy
\widowpenalty=10000
\clubpenalty=10000
\displaywidowpenalty=10000
\begin{tabularx}{\linewidth}{@{}>{\bfseries}p{0.22\linewidth}X@{}}
組織 & OpenAI \\
種別 & モデル \\
時系列 & — \\
\end{tabularx}
\smallskip
{\bfseries 一次資料から整理したtechnical points}
\begin{itemize}[leftmargin=1.5em,itemsep=0.35em]
\item \textbf{一次情報で確認できる事実}: OpenAIの選定済み一次資料では、GPT-3は175B parameterのautoregressive language modelとして、zero-shot・one-shot・few-shotをtask textとdemonstrationだけで指定し、評価時にgradient updateやtask-specific fine-tuningを行わない設定で検証された。 論文はtranslation・question answeringなどの結果と同時に、few-shot learningが苦手なdatasetやlarge web-corpus trainingに起因するmethodological issueも記録しており、few-shot性能を全面的なtask generalityとはみなさない。 paperとOpenAI first-party publicationはいずれも2020-05-28の公開境界として保持し、研究結果とfirst-party publicationを別sourceとして追跡する。
\end{itemize}
\begin{minipage}{\linewidth}
% reader-facing Technical Notes generic boundary/source tail group
\begin{samepage}
{\bfseries 一次資料}
\begingroup\sloppy
\Urlmuskip=0mu plus 2mu\relax
\begin{itemize}[leftmargin=1.5em,itemsep=0.25em]
\item {\scriptsize\url{http://arxiv.org/abs/2005.14165v4}}
\item {\scriptsize\url{https://openai.com/index/language-models-are-few-shot-learners}}
\end{itemize}
\endgroup
\end{samepage}
\end{minipage}
\endgroup
\end{technicalnote}
\medskip
